\chapter{有界线性算子}
\section{有界线性算子与有界线性泛函}
\begin{definition}[线性算子---保证可以拆括号提取系数也就是线性]
设 $X, X_1$ 是数域 $K$ 上的赋范空间，$D(T)$ 是 $X$ 的线性子空间。
映射
\[
T: D(T) \to X_1
\]
称为从 $D(T)$ 到 $X_1$ 的\textbf{线性算子}，若对任意 $x,y \in D(T)$ 与 $a \in K$，
有
\[
T(x+y)=Tx+Ty, \qquad T(ax)=aTx.
\]
称 $D(T)$ 为 $T$ 的\textbf{定义域}。  
若 $D(T)=X$，则称 $T$ 为 $X$ 上到 $X_1$ 中的线性算子。
\end{definition}
\begin{example}[线性算子的具体例子]
设 $X=C[0,1]$ 为定义在区间 $[0,1]$ 上的所有连续实函数所成的线性空间，
$X_1=\mathbb{R}$。

定义映射
\[
T: C[0,1] \to \mathbb{R}, \qquad
T(f)=\int_0^1 f(x)\,dx.
\]

则对任意 $f,g\in C[0,1]$ 与 $a\in\mathbb{R}$，
有
\[
T(f+g)=\int_0^1 (f+g)(x)\,dx
= \int_0^1 f(x)\,dx + \int_0^1 g(x)\,dx = T(f)+T(g),
\]
\[
T(af)=\int_0^1 af(x)\,dx = a\int_0^1 f(x)\,dx = aT(f).
\]
因此 $T$ 是 $C[0,1]$ 上到 $\mathbb{R}$ 中的\textbf{线性算子}。
\end{example}
\begin{note}
   $D(T)$ 是原像集，所以是定义域。  
如果 $D(T)=X$，那么 $T$ 就是 $X$ 到 $X_1$ 中的线性算子；  
否则，$T$ 是从 $D(T)$ 到 $X_1$ 中的线性算子。  

此外，$X, X_1$ 是数域 $K$ 上的赋范空间，  
是指它们的系数在数域 $K$ 上。

\end{note}
\begin{definition}[线性泛函]
设 $X$ 是数域 $K$ 上的赋范空间。  
若 $X$ 上到 $K$ 中的线性算子为线性泛函。

若 $K=\mathbb{R}$，称为\textbf{实线性泛函}；
若 $K=\mathbb{C}$，称为\textbf{复线性泛函}。
\end{definition}
\begin{note}
    线性泛函是特殊的线性算子，他要求映射到标量域中，也就是输出是一个数不是一个向量
\end{note}
\begin{definition}[有界线性算子]
设 $(X,\|\cdot\|)$ 与 $(X_1,\|\cdot\|_1)$ 是赋范空间。
若线性算子 $T:X\to X_1$ 满足：存在一个正常数 $M$，
使得对任意 $x\in X$，都有
\[
\|T(x)\|_1 \le M\|x\|,
\]
则称 $T$ 为从 $X$ 到 $X_1$ 的\textbf{有界线性算子}。
\end{definition}

\begin{definition}[有界线性泛函]
设 $(X,\|\cdot\|)$ 是赋范空间。
若 $X$ 上的线性泛函 $f:X\to K$ 满足：存在一个正常数 $M$，
使得对任意 $x\in X$，都有
\[
|f(x)| \le M\|x\|,
\]
则称 $f$ 为 $X$ 上的\textbf{有界线性泛函}。
\end{definition}
\begin{theorem}[线性算子的前提下一点连续整个联系]
设 $X, X_1$ 是数域 $K$ 上的赋范空间，$T$ 是 $X$ 上到 $X_1$ 中的线性算子。  
如果 $T$ 在某一点 $x_0 \in X$ 连续，则 $T$ 在 $X$ 上连续。
\end{theorem}

\begin{theorem}[连续的必要条件是T是有界线性算子]
设 $X, X_1$ 是数域 $K$ 上的赋范空间，$T$ 是 $X$ 上到 $X_1$ 中的线性算子。  
则 $T$ 连续，当且仅当 $T$ 有界。
\end{theorem}
\begin{definition}[算子的范数---映射到X1上放大X1中元素最大的放大倍数]
设 $T$ 是 $X$ 上到 $X_1$ 中的有界线性算子，定义
\[
\|T\|=\sup_{x\in X,\,x\ne 0}\frac{\|Tx\|_1}{\|x\|},
\]
称 $\|T\|$ 为算子 $T$ 的\textbf{范数}。


\end{definition}
\begin{itemize}
  \item 由有界线性算子的定义可知，算子 $T$ 的范数定义是有意义的。
  
  \item 任意有界线性算子 $T$ 都满足
  \[
  \|Tx\|_1 \le \|T\|\,\|x\|, \qquad (\forall\,x\in X).
  \]
  
  \item 的精确最小上界定义。
  \[
  \|T\|=\inf\{\,M>0 \mid \forall x\in X,\ \|T(x)\|_1\le M\|x\|\,\}.
  \]
  
  \item 这是定义的计算形式，表达算子范数的“极值意义”
  \[
  \|T\|=\sup_{\|x\|=1}\|Tx\|_1
  =\sup_{\|x\|\le1}\|Tx\|_1.
  \]
\end{itemize}
\begin{example}[例 3.1.1]
考虑 $n$ 阶方阵 $(a_{ik})\ (i,k=1,2,\dots,n)$。  
定义映射
\[
A:\mathbb{R}^n \to \mathbb{R}^n,\qquad A x = (a_{ik})x,\ \forall x\in\mathbb{R}^n.
\]
则 $A$ 是 $\mathbb{R}^n$ 上到 $\mathbb{R}^n$ 中的\textbf{有界线性算子}。
\end{example}

\begin{example}[例 3.1.2]
设 $y_0(t)\in C[a,b]$。定义泛函
\[
f:C[a,b]\to \mathbb{R},\qquad
f(x)=\int_a^b x(t)y_0(t)\,dt,\quad \forall x\in C[a,b].
\]
则 $f$ 是 $C[a,b]$ 上的\textbf{有界线性泛函}。
\end{example}

\begin{example}[例 3.1.3]
给定无穷矩阵 $(a_{ik})$ 满足条件
\[
\sum_{i=1}^{\infty}\sum_{k=1}^{\infty}|a_{ik}|^q<\infty,\qquad (q>1).
\]
对每个 $x=\{\xi_k\}\in l^p$，定义映射
\[
A x = \{\eta_i\},\qquad \eta_i=\sum_{k=1}^{\infty}a_{ik}\xi_k,\quad i=1,2,\dots
\]
于是映射
\[
A:l^p \longrightarrow l^q
\]
是 $l^p$ 上的\textbf{有界线性算子}，其中 $\frac{1}{p}+\frac{1}{q}=1$。
\end{example}

\begin{example}[例 3.1.4（反例）]
设 $X=C[a,b]$。考虑算子
\[
T:C^1[a,b]\to C[a,b],\qquad T x(t)=x'(t).
\]
则 $T$ 是 $C[a,b]$ 中到 $C[a,b]$ 中的线性算子，
但\textbf{不是有界线性算子}。
\end{example}
\begin{definition}[有界线性算子空间]\label{def:bounded-operator-space}
设 $X, X_1$ 是数域 $K$ 上的赋范空间，记
\[
\mathcal{B}(X, X_1)
\]
为所有从 $X$ 到 $X_1$ 的有界线性算子全体。  
$\mathcal{B}(X, X_1)$ 称为\textbf{线性赋范空间}，其线性运算与范数分别定义如下：  
对于任意 $A,B\in\mathcal{B}(X, X_1)$ 及 $a\in K$，有
\[
(A+B)(x)=A(x)+B(x),\qquad (aA)(x)=aA(x),
\]
并定义算子范数为
\[
\|A\|=\sup_{x\in X,\,x\ne0}\frac{\|Ax\|}{\|x\|}.
\]
特别地，当 $X=X_1$ 时，常将 $\mathcal{B}(X,X_1)$ 简记为 $\mathcal{B}(X)$。
\end{definition}
% 等价刻画：算子范数收敛 ⇔ 在单位球面上一致收敛
\begin{itemize}
  \item 在 $B(X,X_1)$ 中按范数收敛等价于算子列在 $X$ 的单位球面一致收敛。
  
  \item 即：对于任意的 $A,A_n\in B(X,X_1)$，
  \[
    \|A_n-A\|\xrightarrow[n\to\infty]{}0
    \quad\Longleftrightarrow\quad
    \forall\,\varepsilon>0\ \exists\,N\ \text{s.t.}\ n>N \Rightarrow
    \forall\,x\in S=\{x\in X:\|x\|=1\},\ 
    \|A_nx-Ax\|_{X_1}<\varepsilon .
  \]
  
  \item 在 $B(X,X_1)$ 中按范数收敛也称为一致收敛。
  
  \item 结论中单位球面可替换成任意有界集。
\end{itemize}
\begin{theorem}[对称正定矩阵的算子范数与特征值的关系]
设 $A$ 为 $n$ 阶实矩阵，定义 $B = A^\top A$。  
则 $B$ 为对称非负定矩阵，且
\[
\|A\| = \sqrt{\max_{1\le i\le n}\lambda_i},
\]
其中 $\lambda_i$ 为 $B$ 的特征值。
\end{theorem}

\begin{proof}
由于 $B=A^\top A$ 是实对称矩阵，存在正交矩阵 $P$，使得
\[
B = P^\top 
\begin{pmatrix}
\lambda_1 & & 0\\
& \ddots & \\
0 & & \lambda_n
\end{pmatrix}
P,
\qquad
\lambda_i \ge 0.
\]

令
\[
\xi = P x =
\begin{pmatrix}
\xi_1 \\ \xi_2 \\ \vdots \\ \xi_n
\end{pmatrix},
\]
因为 $P$ 是正交矩阵，所以 $\|P x\| = \|x\|$。

于是有
\[
\|A x\|^2 = x^\top B x 
= x^\top P^\top 
\begin{pmatrix}
\lambda_1 & & 0\\
& \ddots & \\
0 & & \lambda_n
\end{pmatrix}
P x
= \xi^\top
\begin{pmatrix}
\lambda_1 & & 0\\
& \ddots & \\
0 & & \lambda_n
\end{pmatrix}
\xi.
\]

展开可得
\[
\|A x\|^2 
= \lambda_1 \xi_1^2 + \lambda_2 \xi_2^2 + \cdots + \lambda_n \xi_n^2.
\]

由于
\[
\|x\|^2 = \|P x\|^2 = \xi_1^2 + \xi_2^2 + \cdots + \xi_n^2,
\]
记 $S = \{x\in\mathbb{R}^n : \|x\|=1\}$，则对任意 $x\in S$，
\[
\|A x\|^2 = \sum_{i=1}^n \lambda_i \xi_i^2
\le \left( \max_{1\le i\le n}\lambda_i \right)
      \sum_{i=1}^n \xi_i^2
= \left( \max_{1\le i\le n}\lambda_i \right) \|x\|^2.
\]
因此
\[
\|A\|^2 = \sup_{\|x\|=1}\|A x\|^2
\le \max_{1\le i\le n}\lambda_i.
\]

当取 $x_0$ 为对应最大特征值 $\lambda_{\max}$ 的特征向量时，
\[
B x_0 = \lambda_{\max} x_0, \quad \|x_0\|=1,
\]
于是
\[
\|A x_0\|^2 = x_0^\top B x_0 = \lambda_{\max}\|x_0\|^2 = \lambda_{\max}.
\]
由此达到上界，得
\[
\boxed{\|A\|^2 = \max_{1\le i\le n}\lambda_i}, \quad 即\quad
\boxed{\|A\| = \sqrt{\max_i \lambda_i}}.
\]
\end{proof}

\begin{remark}
上式表明：矩阵 $A$ 的算子范数等于 $A^\top A$ 最大特征值的平方根。  
几何意义上，这表示 $A$ 将单位球面映射成一个椭球体，其主轴长度的最大值即为算子范数。
\end{remark}
\begin{theorem}[Banch空间和所有赋范空间组成的空间都是banch空间]
设 $X$ 是赋范空间，$X_1$ 是 Banach 空间，则 $\mathcal{B}(X, X_1)$ 是 Banach 空间。
\end{theorem}
若目标空间 $X_1$ 是完备的，则所有从 $X$ 到 $X_1$ 的有界线性算子所构成的空间 $\mathcal{B}(X,X_1)$ 在算子范数下也是完备的。
\begin{example}
    \[
\mathcal{B}(X, X_1) \text{ 就是所有 } m \times n \text{ 的矩阵组成的空间：}
\]

\[
\mathcal{B}(\mathbb{R}^n, \mathbb{R}^m) = \mathbb{R}^{m \times n}
\]

\[
\text{有限维空间必完备} \;\Rightarrow\; \mathcal{B}(\mathbb{R}^n, \mathbb{R}^m) \text{ 是 Banach。}
\]
\end{example}
\begin{definition}[算子列的逐点收敛（强收敛）]
设 $A, A_n \in \mathcal{B}(X, X_1)$。若对任意 $x \in X$，
\[
A_n x \to Ax \qquad (n \to \infty),
\]
则称算子列 $\{A_n\}$ 逐点收敛于 $A$，或称 $\{A_n\}$ 强收敛于 $A$。
\end{definition}
\textbf{● 若 $\{A_n\}$ 一致收敛于 $A$，则 $\{A_n\}$ 逐点收敛于 $A$。反之不成立。}

\begin{example}
在空间 $\ell^p\,(p \ge 1)$ 中定义算子列 $\{T_n\}$，
\[
T_n x = x_n \qquad (n = 1,2,\dots),
\]
其中 $x = \{\xi_k\} \in \ell^p,\quad x_n = \{\xi_n, \xi_{n+1},\dots\}$。
\end{example}
\section{Banach-Steinhaus（施坦因豪斯）
            定理及其某些应用
}
\begin{theorem}[ (Banach–Steinhaus)一致有界原理]
设 $\{T_a\}(a \in I)$ 是 Banach 空间 $X$ 上到赋范空间 $X_1$ 中的有界线性算子族，
如果对于每个 $x \in X$，
\[
\sup_{a \in I} \|T_a x\| < \infty,
\]
则 $\{\|T_a\|\}(a \in I)$ 是有界集。
\end{theorem}

对于所有 $x \in X$，若集合 $\{\|T_a(x)\| : a \in I\}$ 有界，则
\[
\{\|T_a(x)\| : \|x\| \le 1, \, a \in I\} \text{ 有界}.
\]

（称为一致有界）
\begin{itemize}
  \item 如果一族有界线性算子对每个向量作用出来的值都不会“失控”，即
    \[
    \sup_{a \in I} \|T_a x\| < \infty \qquad (\forall x \in X),
    \]
    那么这些算子点上作用的结果都是有界的。

  \item 于是可以推出：这些算子的算子范数整体也不会“失控”，即存在同一个常数 $M$ 使得
    \[
    \sup_{a \in I} \|T_a\| < \infty.
    \]
    换句话说，这些算子不会突然在某个输入下出现“爆炸”行为。
\end{itemize}
\begin{definition}[疏集]
设 $(X,d)$ 为度量空间，$E \subset X$。若
\[
\operatorname{int}(\overline{E}) = \varnothing,
\]
则称 $E$ 为疏集（nowhere dense）。
\end{definition}
\begin{theorem}[如果算子列一致有界 + 在稠密子集上逐点收敛 → 则在整个空间上强收敛。]
设 $\{T_n\}(n \in \mathbb{N})$ 是赋范空间 $X$ 上到 Banach 空间 $X_1$ 中的有界线性算子列，若满足：
\begin{itemize}
    \item $\{\|T_n\|\}$ 有界；
    \item 对某个稠密子集 $G \subset X$ 中的每个 $y$，都有 $T_n y$ 收敛；
\end{itemize}
则 $T_n$ 强收敛于某个有界线性算子 $T$，并且
\[
\|T\| \le \liminf_{n \to \infty} \|T_n\|.
\]
\end{theorem}
极限算子的大小不会超过原算子序列最终的最小趋势。
\begin{definition}[算子列的逐点收敛（强收敛）]
设 $A, A_n \in \mathcal{B}(X, X_1)$。如果对所有 $x \in X$，
\[
A_n x \to A x \qquad (n \to \infty),
\]
则称 $\{A_n\}$ 逐点收敛于 $A$，或称 $\{A_n\}$ 强收敛于 $A$。
\end{definition}
\begin{theorem}[强收敛完备]
设 $X, X_1$ 是 Banach 空间，则有界线性算子空间 $\mathcal{B}(X, X_1)$
在强收敛意义下完备。
\end{theorem}

\begin{theorem}[强收敛条件]
设 $\{T_n\}(n \in \mathbb{N})$ 是 Banach 空间 $X$ 上到 Banach 空间 $X_1$ 中
的有界线性算子列，则 $T_n$ 强收敛于某个有界线性算子 $T$ 的充要条件是：
\begin{itemize}
    \item $\{\|T_n\|\}$ 有界；
    \item 对某个稠密子集 $G \subset X$ 中的任意 $y$，都有 $T_n y$ 收敛。
\end{itemize}
\end{theorem}
\begin{example}[机械求积公式的收敛性]
在积分的近似计算中，我们通常考虑如下求积公式：
\[
\int_a^b x(t)\, dt \approx \sum_{k=0}^{n} A_k \, x(t_k), 
\qquad (a \le t_0 < t_1 < \cdots < t_n \le b).
\]

例如矩形公式、梯形公式都属于该类求积方法。由于只用一个公式不能保证足够的精确度，我们考虑一组求积公式：
\[
\int_a^b x(t)\, dt \approx \sum_{k=0}^{n} A_k^{(n)} \, x\!\left(t_k^{(n)}\right),
\tag{3.2.1}
\]
其中
\[
a \le t_0^{(n)} < t_1^{(n)} < \cdots < t_n^{(n)} \le b, 
\qquad n = 0,1,2,\ldots
\]
\end{example}


\begin{theorem}[机械求积公式的收敛性问题]
利用公式 (3.2.1) 进行近似计算时，在某些条件下，当 $n \to \infty$ 时误差趋于零。

机械求积公式 (3.2.1) 对每一个连续函数 $x \in C[a,b]$ 都收敛，即
\[
\sum_{k=0}^{n} A_k^{(n)} x\bigl(t_k^{(n)}\bigr)
\;\longrightarrow\;
\int_a^b x(t)\, dt,
\]
当且仅当以下两个条件成立：

1. **存在常数 $M$**，使得
\[
\sum_{k=0}^n \left|A_k^{(n)}\right| \le M, \qquad (n = 0,1,2,\ldots)
\]

2. **公式 (3.2.1) 对每个多项式是收敛的。**

\end{theorem}
\section{开映射定理与闭图像定理}
\begin{definition}[算子乘法--复合函数]
设 $X, X_1, X_2$ 是数域上的赋范空间，$T_1 \in B(X, X_1)$，$T_2 \in B(X_1, X_2)$。
定义算子乘法 $T = T_2 T_1 : X \to X_2$ 为
\[
T x = T_2(T_1 x), \qquad \forall x \in X.
\]
\end{definition}

\begin{theorem}[算子乘法的性质]
对算子乘法 $T = T_2 T_1$，有以下性质成立：

\begin{itemize}
    \item[(1)] 有界性：$T \in B(X, X_2)$，并且
    \[
    \|T\| = \|T_2 T_1\| \le \|T_2\| \, \|T_1\|.
    \]

    \item[(2)] 结合律：
    \[
    (T_4 T_3) T_1 = T_4 (T_3 T_1),
    \]
    其中 $T_1, T_2 \in B(X, X_1)$，$T_3 \in B(X_1, X_2)$，$T_4 \in B(X_2, X_3)$。

    \item[(3)] 分配律：
    \[
    T_3 (T_2 + T_1) = T_3 T_2 + T_3 T_1.
    \]

    \item[(4)] 交换律一般不成立。
\end{itemize}
\end{theorem}
\begin{definition}[算子的逆运算]
设 $X, X_1$ 是线性空间，$T : X \to X_1$ 是线性算子。
如果存在算子 $T_1 : X_1 \to X$ 使得
\[
T_1 T = I_X, \qquad T T_1 = I_{X_1},
\]
则称算子 $T$ 有逆算子（或 $T$ 可逆），$T_1$ 称为 $T$ 的逆算子，并记为 $T_1 = T^{-1}$。
\end{definition}

\begin{theorem}[逆算子的性质]
对于可逆算子 $T$，其逆算子满足以下性质：
\begin{itemize}
    \item[(1)] 线性算子的逆算子也是线性算子；
    \item[(2)] $\left(T^{-1}\right)^{-1} = T$；
    \item[(3)] 若 $T$ 存在逆算子 $T^{-1}$，则 $T$ 是双射；
    \item[(4)] 逆算子存在则必唯一。
\end{itemize}

\end{theorem}
\begin{note}
    和矩阵很像
\end{note}
\begin{definition}[恒等算子与逆算子]
设 $X$ 为线性空间，定义恒等算子（identity operator）
\[
I_X : X \to X, \qquad I_X(x) = x \quad \forall x \in X,
\]
即 $I_X$ 将每一个元素映射为其自身。

设 $T : X \to X_1$ 为线性算子。如果存在线性算子 $T_1 : X_1 \to X$，使得
\[
T_1 T = I_X, \qquad T T_1 = I_{X_1},
\]
则称算子 $T$ 可逆（invertible），并称 $T_1$ 为 $T$ 的逆算子，记作 $T^{-1}$。
\end{definition}
\begin{theorem}[有界逆算子]
设 $T$ 是赋范空间 $X$ 上到赋范空间 $X_1$ 上的线性算子，且存在常数 $m>0$ 使得
\[
\|Tx\| \ge m \|x\| \qquad (x \in X),
\]
则 $T$ 有有界逆算子 $T^{-1}$。
\end{theorem}

\begin{theorem}[算子范数小于1的性质]
设 $X$ 是 Banach 空间，$T \in B(X)$。如果 $\|T\| < 1$，则算子 $I - T$ 有逆算子，且
\[
\|(I - T)^{-1}\| \le \frac{1}{1 - \|T\|}.
\]
\end{theorem}
\newpage
\begin{corollary}
设 $X$ 是 Banach 空间，$T \in B(X)$ 有有界逆算子，则对于任意 $\Delta T \in B(X)$，当
\[
\|\Delta T\| < \frac{1}{\|T^{-1}\|}
\]
时，算子 $S = T + \Delta T$ 有有界逆算子，并且
\[
S^{-1} = \sum_{k = 0}^{\infty} (-1)^k \left( T^{-1} \Delta T \right)^k T^{-1}.
\]
\end{corollary}[扰动自购小就还有逆]
\begin{note}
    一个有逆算子 $T$ 经过足够小的扰动 $\Delta T$ 后，新算子 $S = T + \Delta T$ 仍然有逆，而且逆算子可以写成一个无穷级数。

推论表述：

若 $X$ 是 Banach 空间，$T \in B(X)$ 有有界逆算子 $T^{-1}$，
只要扰动满足
\[
\|\Delta T\| < \frac{1}{\|T^{-1}\|},
\]
则新算子
\[
S = T + \Delta T
\]
仍然是可逆的（即有逆算子 $S^{-1}$），并且它的逆可以展开为 Neumann 级数：
\[
S^{-1} = \sum_{k=0}^{\infty} (-1)^k (T^{-1} \Delta T)^k T^{-1}.
\]
\end{note}

\begin{definition}[有限维空间上的谱]
设涉及到谱问题处，空间皆是定义在复数域上的线性空间。

设 $T$ 是 $n$ 维空间上的线性算子。如果 
\[
Tx = \lambda x
\]
有非零解，称数 $\lambda$ 是算子 $T$ 的特征值。

\begin{itemize}
    \item 所有特征值全体称为算子 $T$ 的谱。
          此时方程 $Tx=\lambda x$ 有非零解，$(T-\lambda I)^{-1}$ 不存在。
    \item 其他 $\lambda$ 称为 $T$ 的正则点。
          此时 $(T-\lambda I)^{-1}$ 存在，且为定义在全空间上的有界线性算子。
    \item 无限维空间中:$(T-\lambda I)^{-1}$ 存在，但不是定义在全空间上，也可能不是有界线性算子。
\end{itemize}

\end{definition}






\begin{note}
  谱是所有特征值，不是特征值的点就是正则点只是这个变量也叫$\lambda  $叫正则点是因为逆存在。\\
  谱：spectrum 
  特征值：eigenvalue
  正则点：regular point 或 resolvent point
\end{note}
\begin{definition}[谱与正则值]
设 $T$ 是 Banach 空间上的有界线性算子。若算子 $(T - \lambda I)^{-1}$ 存在且定义在全空间 $X$ 上，则称数 $\lambda$ 为算子 $T$ 的正则值。此时称
\[
R_\lambda = (T - \lambda I)^{-1}
\]
为算子 $T$ 的预解式。

称所有其他 $\lambda$ 为算子 $T$ 的谱点，算子 $T$ 的谱点全体称为算子 $T$ 的谱，记为 $\sigma(T)$。
\end{definition}

\noindent 算子 $\sigma(T)$ 可分为：

\begin{itemize}
    \item \textbf{点谱：} 若 $T - \lambda I$ 不是一对一的，即方程 $Tx = \lambda x$ 有非零解。
    \item \textbf{连续谱：} 若 $T - \lambda I$ 的值域不是全空间。
\end{itemize}
\begin{theorem}[完备空间上的有界线性算子的谱是有界闭集]
设 $X$ 是 Banach 空间，$T \in B(X)$，则 $\sigma(T)$ 是有界闭集。
\end{theorem}

\begin{corollary}[逆算子判定]
设 $X$ 是 Banach 空间，$T \in B(X)$ 有有界逆算子，则对任意 $\Delta T \in B(X)$，当
\[
\|\Delta T\| < \frac{1}{\|T^{-1}\|}
\]
时，算子 $S = T + \Delta T$ 有有界逆算子，且
\[
S^{-1} = \sum_{k=0}^{\infty} (-1)^k (T^{-1} \Delta T)^k T^{-1}.
\]
\end{corollary}

\begin{theorem}[完备空间算子范数小于1I-T一定有逆算子]
设 $X$ 是 Banach 空间，$T \in B(X)$。如果 $\|T\| < 1$，则算子 $I - T$ 有逆算子，且
\[
\|(I - T)^{-1}\| \le \frac{1}{1 - \|T\| }.
\]
\end{theorem}
\subsection{开映射定理}
\begin{theorem}[开映射定理-完备+有界线性算子=开映射]
设 $X, X_1$ 是 Banach 空间，$T$ 是 $X$ 上到 $X_1$ 上的有界线性算子，则 $T$ 是开映射。
\end{theorem}
有界线性算子就算开映射
\begin{itemize}

  \item \textbf{第 1 步：} 若存在 $\delta>0$，使得
  \[
      S_{X_1}(\theta,\delta) \subset T(S_X(\theta,1)),
  \]
  则 $T$ 是一个开映射。

  \item \textbf{第 2 步：} 若存在 $\delta>0$，使得
  \[
      S_{X_1}(\theta,2\delta) \subset \overline{T(S_X(\theta,1))},
  \]
  则
  \[
      S_{X_1}(\theta,\delta) \subset T(S_X(\theta,1)).
  \]

  \item \textbf{第 3 步：} 存在 $\delta>0$，使得
  \[
      S_{X_1}(\theta,2\delta) \subset \overline{T(S_X(\theta,1))}.
  \]

\end{itemize}
\begin{definition}[开映射——开集的像还是开集]
设 $T$ 是距离空间 $X$ 上到距离空间 $X_1$ 中的一个映射。若对于 $X$ 内任意的开子集 $E$，其像集 $T(E)$ 是 $X_1$ 内的开集，则称 $T$ 是一个开映射。
\end{definition}
开集的像还是开集
\begin{theorem}[Banach 逆算子定理]
设 $X, X_1$ 是 Banach 空间，$T$ 是 $X$ 上到 $X_1$ 上的一一的映射，则 $T$ 的逆算子 $T^{-1}$ 是有界线性算子。
\end{theorem}
一一映射的逆算子还是有界线性算子
\begin{corollary}[都是完备空间，且有常数限制他们两的大小关系就是等价]
设线性空间 $X$ 上的两个范数 $\|\cdot\|_1$ 与 $\|\cdot\|_2$ 都使 $X$ 成为 Banach 空间，
并且存在常数 $C>0$ 使得
\[
    \|x\|_2 \le C \|x\|_1 \qquad (x\in X),
\]
则这两个范数是等价的。

\begin{itemize}
    \item Banach 空间中的范数能够比较即等价。
    \item Banach 空间中的恒等算子 $(X,\|\cdot\|_1)\to (X,\|\cdot\|_2)$ 不一定连续。
\end{itemize}
\end{corollary}
\subsection{闭图像}
\begin{definition}[算子的图像]
设 $X, X_1$ 是赋范空间，$T$ 是从 $X$ 到 $X_1$ 的线性算子，定义其图像为
\[
G(T)=\{(x,Tx)\in X\times X_1: x\in D(T)\},
\]
其中 $D(T)$ 为算子 $T$ 的定义域。
\end{definition}

\begin{definition}[闭算子]
若算子 $T$ 的图像 $G(T)$ 在乘积赋范空间 $X\times X_1$ 中是闭集，则称算子 $T$ 是\textbf{闭算子}。
\end{definition}

\begin{example}
设 $X=X_1=\mathbb{R}$，赋予通常的绝对值范数。定义线性算子
\[
T:\mathbb{R}\to\mathbb{R},\qquad T(x)=2x.
\]
则算子 $T$ 的定义域为 $D(T)=\mathbb{R}$，其图像为
\[
G(T)=\{(x,Tx):x\in\mathbb{R}\}
=\{(x,2x):x\in\mathbb{R}\}.
\]

几何上，$G(T)$ 是 $\mathbb{R}^2$ 中一条直线，它由所有点 $(x,2x)$ 组成：
- 横坐标是输入 $x$，
- 纵坐标是输出 $Tx=2x$。

注意到 $G(T)$ 在 $\mathbb{R}^2$ 中是闭集（因为它是一条通过原点的直线，而直线是闭集），
因此算子 $T$ 是一个闭算子。
也就是说：若 $(x_n,Tx_n)\to (x,y)$，则必有 $y=Tx=2x$。
\end{example}
\begin{theorem}[闭算子 = 图像在极限下封闭：输入与输出都收敛时，极限点仍在图像里。]
设 $X, X_1$ 是赋范空间，$T$ 是 $X$ 中到 $X_1$ 中的线性算子。  
$T$ 是闭算子，当且仅当：对任意 $\{x_n\}\subset D(T)$，如果满足
\[
\lim_{n\to\infty} x_n = x \in X
\quad\text{且}\quad
\lim_{n\to\infty} Tx_n = y \in X_1,
\]
则有 $x\in D(T)$，且 $Tx=y$。
\end{theorem}
\section{Hahn-Banach定理及其推论}
\begin{definition}[半序]
设 $F$ 是一个非空集合，$\preceq$ 是 $F$ 上的一个二元关系。若对任意 $a,\beta,\gamma\in F$，均满足：

\begin{itemize}
  \item[(1)] 自反性：$a\preceq a$；
  \item[(2)] 传递性：$a\preceq \beta,\; \beta\preceq \gamma \;\Rightarrow\; a\preceq \gamma$；
  \item[(3)] 反对称性：$a\preceq \beta,\; \beta\preceq a \;\Rightarrow\; a=\beta$，
\end{itemize}

则称 $\preceq$ 为 $F$ 上的一个\textbf{半序}，并称 $(F,\preceq)$ 为一个\textbf{半序集}。
\end{definition}
\begin{example}
设 $F=\{\{1\},\{1,2\},\{1,2,3\}\}$，在 $F$ 上定义关系
\[
A \preceq B \quad \Longleftrightarrow\quad A\subseteq B ,
\]
即用“集合包含关系”作为比较方式。

那么 $\preceq$ 是 $F$ 上的一个半序，因为：

\begin{itemize}
  \item 自反性：任意集合 $A$ 满足 $A\subseteq A$；
  \item 传递性：若 $A\subseteq B$ 且 $B\subseteq C$，则 $A\subseteq C$；
  \item 反对称性：若 $A\subseteq B$ 且 $B\subseteq A$，则 $A=B$。
\end{itemize}

因此 $(F,\preceq)$ 构成一个半序集。
\end{example}
\begin{definition}[全序]
若在半序集 $(F,\preceq)$ 中，对任意 $a,\beta\in F$，总有 $a\preceq \beta$ 或 $\beta\preceq a$ 至少一个成立，则称 $(F,\preceq)$ 为\textbf{全序集}。
\end{definition}
\begin{example}
设 $F=\mathbb{Z}$ 为全体整数集合，定义二元关系
\[
a \preccurlyeq b \quad \Longleftrightarrow \quad a \le b,
\]
其中 $\le$ 为通常的大小关系。

显然，对任意 $a,b\in\mathbb{Z}$，必有 $a\le b$ 或 $b\le a$，因此
\[
(\mathbb{Z},\preccurlyeq)
\]
是一个全序集。

直观地说，全体整数按照从小到大的自然排列方式，本身构成一个典型的全序结构。
\end{example}

\begin{definition}[极大元]
设 $(A,\leq)$ 是一个半序集，$\beta\in A$。如果对任意 $a\in A$，只要有
\[
\beta\leq a,
\]
就必有 $a=\beta$，则称 $\beta$ 是 $(A,\leq)$ 的一个极大元。
（直观地说：没有比 $\beta$ 更大的元素。）
\end{definition}

\begin{definition}[上界]
设 $B\subset A$ 是一个非空子集。若对任意 $a\in B$ 都有
\[
a\leq \beta,
\]
则称 $\beta$ 是 $B$ 的一个上界。
（直观地说：$\beta$ 比 $B$ 中所有元素都大。）
\end{definition}

在极大元与上界存在的前提下（它们本身不一定存在），有以下说明：
\begin{itemize}
  \item 极大元与上界一般都不唯一；
  \item 集合 $B$ 的极大元一定属于 $B$，而 $B$ 的上界不一定属于 $B$。
\end{itemize}

\begin{theorem}[Zorn 引理]\label{thm:Zorn}
设 $(A,\leq)$ 是一个半序集。如果 $(A,\leq)$ 中任意全序子集（链）在 $A$ 中都有上界，
则 $(A,\leq)$ 中必存在极大元。
\end{theorem}
\begin{theorem}[实空间的 Hahn--Banach 定理（受控延拓定理）]\label{thm:HB-real}
设 $M$ 是实线性空间 $X$ 的一个线性子空间，$p:X\to\mathbb{R}$ 满足  
对任意 $x,y\in X$、$a\ge0$，
\[
p(x+y)\le p(x)+p(y),\qquad p(ax)=a\,p(x).
\]
设 $f$ 是 $M$ 上的线性泛函，且满足
\[
f(x)\le p(x),\qquad \forall x\in M.
\]
则存在 $X$ 上的线性泛函 $F$，使得
\[
F(x)=f(x),\quad \forall x\in M,
\qquad\text{且}\qquad
-p(-x)\le F(x)\le p(x),\quad \forall x\in X.
\]
\end{theorem}
\begin{theorem}[赋范空间的 Hahn--Banach 定理（保范延拓定理）]\label{thm:HB-norm}
设 $G$ 是复赋范空间 $X$ 的线性子空间，$f$ 是 $G$ 上的有界线性泛函。
则 $f$ 可以保范延拓到整个空间 $X$ 上，即存在 $X$ 上的有界线性泛函
$F$，使得
\begin{enumerate}
  \item $F(x)=f(x),\quad \forall x\in G$；
  \item $\|F\|_{X}=\|f\|_{G}$.
\end{enumerate}
\end{theorem}

\begin{corollary}\label{cor:HB-support}
设 $X$ 是赋范空间，则对任意 $x_0\in X$ 且 $x_0\neq0$，必存在
$X$ 上的有界线性泛函 $f$，使得
\[
\|f\|=1,\qquad f(x_0)=\|x_0\|.
\]
\end{corollary}
\begin{corollary}\label{cor:HB-distance}
设 $G$ 是赋范空间 $X$ 的线性子空间，$x_0\in X$。若
\[
d=d(x_0,G):=\inf_{x\in G}\|x-x_0\|>0,
\]
则存在 $X$ 上的有界线性泛函 $f$，使得
\[
f(x)=0,\quad \forall x\in G;\qquad
\|f\|=\frac1d;\qquad
f(x_0)=1.
\]
\end{corollary}

\begin{corollary}\label{cor:HB-norm1}
设 $X$ 是赋范空间，则对任意 $x_0\in X$、$x_0\neq0$，必存在
$X$ 上的有界线性泛函 $f$，使得
\[
\|f\|=1,\qquad f(x_0)=\|x_0\|.
\]
\end{corollary}
\begin{definition}[共轭空间 / 对偶空间]
设 $X$ 是赋范空间，记
\[
X^* = B(X,\mathbb K),
\]
其中 $\mathbb K=\mathbb R$ 或 $\mathbb C$，$B(X,\mathbb K)$ 表示
$X$ 到 $\mathbb K$ 的所有有界线性泛函的全体。称 $X^*$ 为 $X$ 的共轭空间
（或对偶空间）。
\end{definition}

\begin{definition}[共轭算子 / 对偶算子]
设 $X,X_1$ 是赋范空间，$T\in B(X,X_1)$。对任意 $f\in X_1^*$，定义
\[
(T^* f)(x) = f(Tx),\qquad \forall x\in X.
\]
称算子 $T^*:X_1^*\to X^*$ 为 $T$ 的共轭算子（或对偶算子）。
\end{definition}

\begin{proposition}
若 $X$ 是赋范空间，则其共轭空间 $X^*$ 在算子范数下一定是 Banach 空间。
\end{proposition}
\begin{theorem}[有界线性算子的共轭算子的性质]\label{thm:adjoint-properties}
设 $X,X_1$ 为赋范空间，$T\in B(X,X_1)$，记其共轭算子为
$T^*:X_1^*\to X^*$。则 $T^*$ 具有以下性质：
\begin{enumerate}
  \item $T^*$ 是有界线性算子，并且
        \[
        \|T^*\|=\|T\|;
        \]
  \item 对任意 $a\in\mathbb K$，
        \[
        (aT)^* = aT^*;
        \]
  \item 若 $T_1,T_2\in B(X,X_1)$，则
        \[
        (T_1+T_2)^* = T_1^* + T_2^*;
        \]
  \item 若 $T_1\in B(X_1,X_2)$、$T_2\in B(X,X_1)$，则
        \[
        (T_1T_2)^* = T_2^*T_1^*;
        \]
  \item 若 $T$ 有有界逆算子 $T^{-1}\in B(X_1,X)$，则 $T^*$ 也有有界逆算子，
        并且
        \[
        (T^*)^{-1} = (T^{-1})^*.
        \]
\end{enumerate}
\end{theorem}
\begin{theorem}[F.\ Riesz]\label{thm:Riesz-Cab}
设 $f$ 是 $C[a,b]$ 上的有界线性泛函，则存在区间 $[a,b]$ 上的
一个有界变差函数 $v(t)$，使得
\begin{equation}\label{eq:Riesz-representation}
    f(x)=\int_a^b x(t)\,dv(t),\qquad \forall x\in C[a,b],
\end{equation}
并且
\[
    \|f\| = V_a^b(v),
\]
其中 $V_a^b(v)$ 表示 $v$ 在 $[a,b]$ 上的全变差。

反之，$[a,b]$ 上的任意有界变差函数 $v(t)$，由式
\eqref{eq:Riesz-representation} 都定义了 $C[a,b]$ 上的一个有界线性泛函。
\end{theorem}

记
\[
V_0[a,b]
=\bigl\{
v(t)\in V[a,b]\;\bigm|\;
v(t+0)=v(t),\ \forall\,a<t<b;\ v(a)=0
\bigr\},
\]
则有
\[
(C[a,b])^* = V_0[a,b].
\]
\begin{theorem}\label{thm:Lp-dual}
设 $1<p<\infty$，$f$ 是 $L^p[a,b]$ 上的有界线性泛函，则存在唯一的
$y\in L^q[a,b]$（其中 $1/p+1/q=1$），使得
\begin{equation}\label{eq:Lp-repr}
    f(x)=\int_a^b x(t)\,y(t)\,dt,\qquad \forall x\in L^p[a,b],
\end{equation}
并且
\[
    \|f\|=\|y\|
    =\biggl(\int_a^b |y(t)|^q\,dt\biggr)^{1/q}.
\]
反之，对任意 $y\in L^q[a,b]$，由式 \eqref{eq:Lp-repr} 都定义了
$L^p[a,b]$ 上的一个有界线性泛函。
\end{theorem}

由此得到对偶空间的表述：
\[
\bigl(L^p[a,b]\bigr)^* = L^q[a,b],\quad \forall\,1<p<\infty;
\qquad
\bigl(L^1[a,b]\bigr)^* = L^\infty[a,b].
\]
\begin{theorem}\label{thm:c-dual}
设 $c$ 为所有收敛数列构成的赋范空间，$f\in c^*$。
则存在常数 $a$ 及一列 $\{a_n\}\in \ell^1$，使得
对每个 $x\in c$，若记 $x=\{\xi_n\}$，都有
\begin{equation}\label{eq:c-dual-repr}
    f(x)
    = a\lim_{n\to\infty}\xi_n
      + \sum_{n=1}^{\infty} a_n \xi_n,
\end{equation}
并且
\[
    \|f\| = |a| + \sum_{n=1}^{\infty}|a_n|.
\]

反之，若给定常数 $a$ 及一列 $\{a_n\}\in \ell^1$，
则式 \eqref{eq:c-dual-repr} 定义了 $c$ 上的一个有界线性泛函。

特别地，有
\[
    c^* = \ell^1.
\]
\end{theorem}
\section*{一、自反性}

\begin{definition}[二次共轭空间]
设 $X$ 是赋范空间，$X^*$ 是它的共轭空间。
记
\[
X^{**}=(X^*)^*,
\]
称 $X^{**}$ 为 $X$ 的二次共轭空间。
\end{definition}

\begin{definition}[典型映射（典型嵌入）]
设 $x\in X,\ f\in X^*$。定义
\[
F_x(f)=f(x),\qquad \forall f\in X^* .
\]
则 $F_x\in X^{**}$。由此得到映射
\[
T:X\to X^{**},\qquad T(x)=F_x,
\]
称 $T$ 为 $X$ 到其二次共轭空间的典型映射（典型嵌入）。
\end{definition}

\begin{proposition}[典型映射是线性算子]
对任意 $x_1,x_2\in X,\ a\in\mathbb K,\ f\in X^*$，有
\[
T(x_1+x_2)(f)=F_{x_1+x_2}(f)
= f(x_1+x_2)=f(x_1)+f(x_2)
=F_{x_1}(f)+F_{x_2}(f)
= (T(x_1)+T(x_2))(f),
\]
\[
T(ax_1)(f)=F_{ax_1}(f)=f(ax_1)=a f(x_1)=aF_{x_1}(f)=a\,T(x_1)(f).
\]
故 $T(x_1+x_2)=T(x_1)+T(x_2)$，$T(ax_1)=aT(x_1)$，于是 $T:X\to X^{**}$ 为线性算子。
\end{proposition}
\begin{proposition}
典型映射 $T:X\to X^{**},\;x\mapsto F_x$ 是有界线性算子，且对一切 $x\in X$，
\[
\|T(x)\|=\|F_x\|=\|x\|.
\tag{$*$}
\]
\end{proposition}

\begin{proof}
对任意 $x\in X$ 与 $f\in X^*$，有
\[
|F_x(f)|=|f(x)|\le \|f\|\,\|x\|.
\]
取对所有 $\|f\|\le1$ 的上确界，得
\[
\|F_x\|=\sup_{\|f\|\le1}|F_x(f)|
      \le \|x\|.
\tag{1}
\]

另一方面，由 Hahn–Banach 定理的推论可知：对每个 $x\in X$，存在
$f_0\in X^*$ 使得
\[
\|f_0\|=1,\qquad f_0(x)=\|x\|.
\]
于是
\[
\|F_x\|\ge |F_x(f_0)|
         = |f_0(x)|
         = \|x\|.
\tag{2}
\]
由 (1) 与 (2) 得 $\|F_x\|=\|x\|$，即式 $(*)$ 成立，从而 $T$ 为有界线性算子且为等距映射。
\end{proof}

由 $(*)$ 可知 $T:X\to X^{**}$ 是从 $X$ 到 $X^{**}$ 的等距同构到其像 $T(X)$，
故可将 $X$ 等距地看作 $X^{**}$ 的线性子空间。
\begin{definition}[自反空间]
若赋范空间 $X$ 满足
\[
X = X^{**},
\]
则称 $X$ 为自反空间。
\end{definition}

\begin{remark}
若 $X$ 是自反空间，则 $X$ 与 $X^{**}$ 等距同构。
反之，$X$ 与 $X^{**}$ 等距同构时，$X$ 不一定是自反空间（未必真的有 $X=X^{**}$ 的恒等识别）。
\end{remark}

\begin{example}
\begin{itemize}
  \item $\mathbb{R}^n$ 是自反空间；
  \item 对任意 $1<p<\infty$，$L^p[a,b]$ 是自反空间；
  \item 对任意 $1<p<\infty$，$\ell^p$ 是自反空间。
\end{itemize}
\end{example}
\section*{二、弱收敛}

\begin{definition}[弱收敛]\label{def:weak-conv}
设 $X$ 是赋范空间，$x_0,x_n\in X\ (n=1,2,\dots)$。
若对一切 $f\in X^*$ 都有
\[
f(x_n)\xrightarrow[n\to\infty]{} f(x_0),
\]
则称数列 $\{x_n\}$ \emph{弱收敛}于 $x_0$，记为
\[
x_n \rightharpoonup x_0\quad (n\to\infty),
\]
或记为 $x_n \xrightarrow{w} x_0$。此时称 $x_0$ 为 $\{x_n\}$ 的\emph{弱极限}。
\end{definition}

一些基本性质：
\begin{itemize}
  \item 弱收敛的极限唯一；
  \item 若 $\{x_n\}$ 弱收敛，则数列 $\{\|x_n\|\}$ 有界；
  \item 若 $\{x_n\}$ 强收敛于 $x_0$（即 $\|x_n-x_0\|\to0$），
        则 $\{x_n\}$ 必弱收敛于 $x_0$，反之不然。
\end{itemize}